%%%%%%%%%%%%%%%%%%%%%%%%%%%%%%%%%%%%%%%%%%%%%%%%%%%%%%%%%%%%%%%%%%%%%%%%
\section{Open problems}\label{sec:open}
%%%%%%%%%%%%%%%%%%%%%%%%%%%%%%%%%%%%%%%%%%%%%%%%%%%%%%%%%%%%%%%%%%%%%%%%

\begin{enumerate}
\item \emph{A polynomial lower bound.} Prove $M_n\ge\mu^nn^{-\chi}$ for some $\chi$, equivalently
$M_{2m}=O(m^{\chi'}M_m^2)$ for some $\chi'$. This is the main gap. By Proposition~\ref{prop:linrec}, it would follow from a
lower bound, polynomial in $m$, for the probability that at most a fixed number of arches pass over
the middle gap of a meander of order $2m$. It is implied by the existence of an exponent.
\item \emph{Improving the square root.} Prove $L(n)=o(\sqrt n)$. Two routes are stated precisely: short
balanced separators for a not too small fraction of meanders (Proposition~\ref{prop:fraction}), or
the depth hypothesis of Proposition~\ref{prop:depth} with some $\beta>1$. For the second, the
correlation statement for shallow arch systems in Section~\ref{sec:improve} is a cleaner sufficient
condition that can first be tested exactly.
\item \emph{An upper bound beyond $1/n$.} Prove $u_n=O(n^{-\kappa})$ for some $\kappa>1$. Theorem~\ref{thm:threehalves} gives $3/2$ only along a subsequence. A related open
statement is $M_{n+1}\le\mu M_n$ for every $n$, which holds for $n\le31$.
\item \emph{Middle cuts.} Show that cuts with both sides large carry vanishing total probability,
that is, that $(u_n)$ is subexponential, and deduce that a positive fraction of meanders is prime.
\item \emph{Bounded labels, sharp form.} Prove
$M^{[-1,2]}_m\asymp\Lambda^mm^{-\alpha_{\mathrm b}}$, and find the exponents of wider bands of
labels. For every band of labels that contains $0$ and $1$, the corresponding family is closed under
products, so it has a growth rate, but the exponents of the wider bands are not known, nor whether they approach $\alpha_*$ as the
band widens.
\item \emph{The exponent.} Prove that $u_n=n^{-\alpha+o(1)}$ for some $\alpha$, and that
$\alpha=\alpha_*$. Paper~A explains the conjectured value~\cite{PaperA}, and Paper~D studies a
continuum counterpart~\cite{PaperD}.
\end{enumerate}
